\documentclass[11pt]{article}

\usepackage[margin=25mm]{geometry}

\usepackage{amsmath,amssymb,fontspec,unicode-math}

\IfFontExistsTF{DejaVu Serif}{\setmainfont{DejaVu Serif}}{\setmainfont{Latin Modern Roman}}

\IfFontExistsTF{Latin Modern Math}{\setmathfont{Latin Modern Math}}{\setmathfont{DejaVu Math TeX Gyre}}

\usepackage{xurl,hyperref}

\setlength{\parindent}{0pt}\setlength{\parskip}{6pt}

\begin{document}

\title{LittleBit: Ultra Low-Bit Quantization via Latent Factorization}

\author{ELVIA preliminary AI analysis}\date{21 October 2026}\maketitle

Preliminary assessment: Minor.\par

Independent scientific review remains pending. Author adoption has not been requested.

The SVD error comparison omits truncation cross terms. The public attention path does not establish latent-cache implementation, and the bit-budget interpretation remains uncertain.

\section*{Review scope}

Latest-version, claim-focused review of the cited author-paper passages and applicable current author-release paths. The stated counterexamples and scope controls were checked analytically; no author program, old verifier, training run or benchmark was executed.

\section*{Problem 1: Quantization error alone does not order full reconstruction error}

Source: Appendix A, Proposition 2, PDF p.15, Eq.(15); Claim 2 SVD setup

Assessment: supported; impact: minor

\subsection*{Source claim (paraphrase)}

Proposition 2 compares \(\|Q(A)+Q(B)-A-B\|_F\) with \(\|Q(A+B)-A-B\|_F\) and infers the same ordering for approximation of the original matrix \(W\).

\subsection*{Analysis}

Put \(T=W-A-B\), \(d=Q(A)+Q(B)-A-B\), and \(e=Q(A+B)-A-B\). Actual errors are \(\|T-d\|_F\) and \(\|T-e\|_F\), so their squared difference contains \(-2\langle T,d-e\rangle_F\). A complete counterexample uses the same entrywise sign quantizer on each nonzero entry and orthogonal matrices \[U=\frac13\begin{pmatrix}1&-2&-2\\-2&1&-2\\-2&-2&1\end{pmatrix},\qquad V=\frac13\begin{pmatrix}-1&-2&-2\\-2&2&-1\\-2&-1&2\end{pmatrix}.\] Set \(W=U\operatorname{diag}(3,2,1)V^\top\), \(A=3u_1v_1^\top\), \(B=2u_2v_2^\top\). Thus \(A+B\) is genuinely the rank-two truncated SVD with a nonzero third-mode tail. Direct rational arithmetic gives \(\|d\|_F^2=7/3<32/9=\|e\|_F^2\), but total squared errors \(\|W-Q(A)-Q(B)\|_F^2=10/3>3=\|W-Q(A+B)\|_F^2\). No zero-entry convention, numerical SVD or changed quantizer explains the reversal.

\subsection*{Finding and scope}

The stated sufficient condition is incomplete; learned residual quantization-aware training itself is not refuted.

\subsection*{Proposed correction}

The exact criterion is \(\|d\|_F^2-\|e\|_F^2<2\langle T,d-e\rangle_F\). If the cross terms vanish, the simpler comparison becomes sufficient under that additional condition.

An explicit local correction is described; author adoption and a corresponding published erratum are not established.

\section*{Problem 2: Factorized weights do not themselves imply a latent KV cache}

Source: Current arXiv v5 Section 5, PDF p.10; Figure 5 and Table 4, PDF p.9; quantization/modules/attention.py, PhiQKVSplitAttention.forward

Assessment: supported; impact: minor

\subsection*{Source claim (paraphrase)}

Section 5 states that latent vectors are cached rather than full-dimensional keys/values, while Table 4 labels its cache reduction as estimated. The current author quantization/modules/attention.py instead projects full-width keys and values, applies RoPE and passes those tensors to Cache.update.

\subsection*{Analysis}

Weight-factor multiplication and cache representation are distinct. Passing full-width key/value tensors does not implement the proposed latent state by itself. A separately implemented cache could transform them, but no such implementation is established in this traced path. The paper’s conceptual and estimated qualifiers remain material.

\subsection*{Finding and scope}

Latent-cache implementation and measured storage savings are not established by this source path.

\subsection*{Proposed correction}

Identify cache reduction as a prospective/estimated design unless the compatible state layout and decode path are supplied. No new cache algorithm is validated here.

An explicit local correction is described; author adoption and a corresponding published erratum are not established.

\section*{Problem 3: The rank convention in the bit-budget table is unresolved}

Source: Appendix D, PDF pp.19–20, Eqs.(16)–(17); Appendix H, Table 13, PDF p.24

Assessment: needs\_evidence; impact: needs\_evidence

\subsection*{Source claim (paraphrase)}

Appendix D counts two equal-rank paths, while Table 13 reports latent dimension \(R\) and nominal bit-per-weight labels for a primary-path microbenchmark.

\subsection*{Analysis}

Interpreting the table’s rank as the rank of each path approximately doubles the nominal bit budget. A primary-only or total-rank convention explains much of that discrepancy, and the paper explicitly identifies primary-path timing. Kernel configuration and included scaling overhead are needed to choose the correct accounting.

\subsection*{Finding and scope}

The timing table is not disproved; the rank and payload convention needs additional evidence.

\section*{Limits}

Preliminary AI assessment; independent human scientific adjudication remains pending.\par

An analytic counterexample or a release discrepancy does not by itself remeasure the paper’s reported experiments.\par

\section*{Source identity}

arXiv:2506.13771v5 (current author version; updated 2026-02-05T01:59:26Z)

\url{https://arxiv.org/pdf/2506.13771v5}

Reviewed PDF SHA-256: \nolinkurl{54508221ea5b5d77d98012cbd0b4f5f08351523a4a1b58d85f7f32691bf63ea8}

The selected author/publisher source is the same work as the original case, matched by exact title, author identity and official version channel. The exact selected PDF SHA-256 was recomputed in the latest-version pass for Issue \#27. The PDF is a retained author/publisher copy; its original retrieval time is stated only when supported by matching acquisition metadata. File modification time was not substituted.

\end{document}
